% =====================================================================
%  methodology_c78_additions.tex
%  Drop-in blocks for methodology.tex (line numbers at the current
%  GitHub head, 1540 lines). Each block states WHERE it goes. Paste the
%  LaTeX between the markers. Labels defined here: sec:meth-ctrl,
%  sec:meth-lax, sec:meth-cost, sec:meth-c8-space, tab:design-space-c78,
%  tab:ctrl-bounds, tab:lax-study, eq:k-ari, eq:g-ctrl-meth,
%  eq:k-max-derived, eq:lax-def, eq:ktarget-3d, eq:ktarget-c8-fit,
%  eq:refl-budget-c8.
%
%  Style: short paragraphs, no semicolons, no em dashes, every displayed
%  formula followed by a term list with units, enumerations as lists.
% =====================================================================

%% ===================================================================
%% BLOCK M1  (methodology.tex, line 233)
%% Replace "over a five-dimensional design space" by the sentence below.
%% ===================================================================
%% >>> M1
over a design space of four to six variables, depending on the
campaign (Tables~\ref{tab:design-space}, \ref{tab:design-space-c56}
and~\ref{tab:design-space-c78}).
%% <<< M1

%% ===================================================================
%% BLOCK M2  (methodology.tex, line 87, caption of the geometry figure)
%% Replace "Ranges for all six design variables" by
%% "Ranges for the six design variables of Campaigns 5 to 7".
%% ===================================================================

%% ===================================================================
%% BLOCK M3  (methodology.tex, insert after line 324, the end of
%% tab:design-space-c56)
%% ===================================================================
%% >>> M3
Campaigns 7 and 8 revise the box a third time. Campaign 7 lowers the
enrichment cap to \SI{12}{wt\%} on the as-built peripheral ring, so
that the search box becomes 3.0 to \SI{10.43}{wt\%}, as a consequence
of the controllability screen of Section~\ref{sec:meth-ctrl}.
Campaign 8 fixes the pitch, collapses the two enrichment variables to
one, raises the cap to \SI{16}{wt\%} and re-bounds the reflector to the
budget that remains once a core barrel and a downcomer are reserved
(Section~\ref{sec:meth-c8-space}).

\begin{table}[htbp]
  \centering
  \caption{Design-space revisions of Campaigns 7 and 8. Variables not
  listed keep the bounds of Table~\ref{tab:design-space}. Campaign 8
  has four design variables.}
  \label{tab:design-space-c78}
  \begin{tabular}{llcc}
    \toprule
    Variable & Meaning & Bounds & Unit \\
    \midrule
    $e_\mathrm{in},\,e_\mathrm{out}$ & enrichments (C7)
      & 3.0 -- 10.435 & wt\% $^{235}$U \\
    $e$ & single assembly enrichment (C8)
      & 3.0 -- 13.913 & wt\% $^{235}$U \\
    $p$ & pin pitch (C8, fixed) & 1.26 & cm \\
    $t_\mathrm{refl}$ & reflector thickness (C8) & 2.0 -- 5.66 & cm \\
    \bottomrule
  \end{tabular}
\end{table}
%% <<< M3

%% ===================================================================
%% BLOCK M4  (methodology.tex, tab:constraints, lines 331-349)
%% Replace the tabular body by the one below. Adds the control screen
%% and the recorded first-two-banks reading, and updates the ceiling.
%% ===================================================================
%% >>> M4
  \begin{tabular}{lll}
    \toprule
    Constraint & Definition & Evaluation level \\
    \midrule
    $g_{k\min}$       & $1.02 - k_\mathrm{BOL} \le 0$                      & assembly to C5, \textbf{core} from C6 \\
    $g_{k\max}$       & $k_\mathrm{BOL} - k_{\max} \le 0$                  & assembly to C5, \textbf{core} from C6, $k_{\max}$ of Table~\ref{tab:ctrl-bounds} \\
    $g_\mathrm{enr}$  & $e_\mathrm{zoned,max} - e_\mathrm{cap} \le 0$      & algebraic, $e_\mathrm{cap}$ per campaign \\
    $g_\mathrm{peak}$ & $F_{\Delta H} - 2.0 \le 0$                         & assembly in C1 to C3, \textbf{core} from C4 \\
    $g_\mathrm{geom}$ & $R_\mathrm{env}(p) + t_\mathrm{refl} + c_\mathrm{v} - 90 \le 0$ & core, analytic, $c_\mathrm{v}$ of Section~\ref{sec:meth-c8-space} \\
    $g_\mathrm{ctrl}$ & $k_\mathrm{ARI} - 0.99 \le 0$                      & \textbf{core}, all sixteen regulating banks in, from C7 \\
    $g_\mathrm{ctrl,12}$ & $k_\mathrm{RE12} - 0.99$                        & core, first two banks in, from C8, \emph{recorded, not constrained} \\
    \bottomrule
  \end{tabular}
%% <<< M4

%% ===================================================================
%% BLOCK M5  (methodology.tex, insert after line 447, the end of
%% subsection "Reactivity window basis")
%% ===================================================================
%% >>> M5
\subsection{Controllability screen and the reactivity ceiling}
\label{sec:meth-ctrl}

Through Campaign 6 the reactivity ceiling proxied controllable excess
reactivity without a rod model. From Campaign 7 the proxy is replaced
by a measurement. Every evaluation solves the beginning-of-life core a
second time with the control rod assemblies of the sixteen central
assemblies inserted, which are the four regulating banks RE1 to RE4 of
the NuScale-inspired seven-bank scheme adopted in Campaign 6, with the
three shutdown banks of the outer sixteen assemblies withdrawn. The
absorber is B$_4$C.

\begin{equation}
k_\mathrm{ARI}(\mathbf{x}) = k_\mathrm{eff}^\mathrm{core}\big(\mathbf{x},\ \mathrm{RE1\ to\ RE4\ inserted}\big)
\label{eq:k-ari}
\end{equation}

\noindent where the symbols have the following meaning.

\begin{itemize}
  \item $k_\mathrm{ARI}$ is the core effective multiplication factor at
        beginning of life with all sixteen regulating-bank assemblies
        inserted, dimensionless.

  \item $\mathbf{x}$ is the design vector.
\end{itemize}

The screen demands an operating subcriticality margin.

\begin{equation}
g_\mathrm{ctrl}(\mathbf{x}) = k_\mathrm{ARI}(\mathbf{x}) - (1 - \Delta\rho_\mathrm{op}) \le 0
\label{eq:g-ctrl-meth}
\end{equation}

\noindent where the symbols have the following meaning.

\begin{itemize}
  \item $g_\mathrm{ctrl}$ is the constraint value, dimensionless,
        negative when satisfied.

  \item $\Delta\rho_\mathrm{op}$ is the operating margin, \num{0.01} in
        units of $k$, that is \SI{1000}{pcm}.
\end{itemize}

The margin is a modelling choice. An insertion of the regulating banks
alone with a \SI{1000}{pcm} margin is not a licensing shutdown margin,
which additionally requires the shutdown banks, the highest-worth
assembly assumed stuck out, and the cold, xenon-free state. The screen
therefore tests whether the regulating system can hold the fresh core
subcritical at power conditions, which is the question the reactivity
ceiling was standing in for.

Campaign 8 adds a second reading with only the first two banks
inserted, RE1 in the four central assemblies and RE2 in the four
diagonal assemblies of the middle ring. It is recorded as
$g_\mathrm{ctrl,12}$ and is not a constraint, so that the front can be
split afterwards into designs controllable with two banks and designs
needing all four, without shrinking the feasible region in advance.

The reactivity ceiling of $g_{k\max}$ can now be derived from the
measured bank worth rather than adopted.

\begin{equation}
k_{\max} = (1 - \Delta\rho_\mathrm{op}) + \min_{\mathbf{x} \in \mathcal{S}} \rho_\mathrm{bank}(\mathbf{x})
\label{eq:k-max-derived}
\end{equation}

\noindent where the symbols have the following meaning.

\begin{itemize}
  \item $k_{\max}$ is the derived ceiling on the unrodded core
        eigenvalue, dimensionless.

  \item $\rho_\mathrm{bank}$ is the bank worth, the difference between
        the unrodded and the rodded eigenvalues, dimensionless.

  \item $\mathcal{S}$ is the sample of designs on which the worth was
        measured.
\end{itemize}

Table~\ref{tab:ctrl-bounds} lists the ceilings used and derived.

\begin{table}[htbp]
  \centering
  \caption{Reactivity ceilings by campaign and their provenance. The
  three-dimensional values come from the axial study of
  Section~\ref{sec:meth-lax}.}
  \label{tab:ctrl-bounds}
  \begin{tabular}{llll}
    \toprule
    Campaigns & $k_{\max}$ & Basis & Provenance \\
    \midrule
    C1 to C6 & 1.35  & assembly, then core & adopted screening value \\
    retrospective & 1.368 & core, 2D & all seven banks, \SI{1000}{pcm}, eight C6 designs \\
    C7 & 1.126 & core, 2D & four regulating banks, eight C6 designs, worths \num{12174} to \SI{12718}{pcm} \\
    C7 measured & 1.099 to 1.217 & core, 2D & implied by each design's own worth, 60 designs \\
    C8 & 1.166 & core, 2D from 3D & smallest 3D worth \SI{14307}{pcm}, times $L_\mathrm{ax}$ \\
    \bottomrule
  \end{tabular}
\end{table}

Two lessons of Campaign 7 govern the use of the derived ceiling. A
minimum over a small sample is not a robust estimator, and the
Campaign 7 value of 1.126 rejected fourteen designs whose own measured
worth showed them controllable (Section~\ref{sec:res-c7-screens}).
And the derived ceiling can only be a surrogate for a quantity that is
measured on every evaluation anyway. From Campaign 8 the measured
screen is the constraint that does the work, and the derived ceiling
is retained as a screening bound with its provenance stated.

The retrospective value of 1.368 also settles the status of the
ceiling used in Campaigns 1 to 6. Derived from the worth of all 32
control rod assemblies with the same margin, it lies two per cent
above the adopted 1.35, which is therefore the measured all-bank
controllability bound of this core to within that tolerance.
%% <<< M5

%% ===================================================================
%% BLOCK M6  (methodology.tex, insert after line 701, the end of
%% section "Leakage-corrected reactivity target (Route B)")
%% ===================================================================
%% >>> M6
\subsection{Axial leakage factor and the three-dimensional target}
\label{sec:meth-lax}

The Route B target of Equation~\eqref{eq:ktarget} corrects the
infinite-lattice depletion for the radial leakage of the
two-dimensional core. The transport model has no axial extent, so the
axial leakage of the finite core is absent from every cycle length of
Campaigns 1 to 7. A dedicated study measured it before Campaign 8.

The three-dimensional model extends the two-dimensional core to a
finite active height of \SI{120}{cm}, adds a \SI{15}{cm} slab of
borated water above and below the fuel as the axial reflector, and
places a vacuum boundary beyond the slabs. The radial reflector spans
the full height. The materials and the zoned enrichment map are those
of the campaign evaluator, so the only difference between the two
models is the axial extent. Each paired solve ran at \num{150000}
particles and 200 batches with 80 inactive, the raised inactive count
allowing the axial source shape to converge, in two seeds.

The axial leakage factor is the ratio of the two eigenvalues.

\begin{equation}
L_\mathrm{ax}(\mathbf{x}) = \frac{k_\mathrm{eff}^\mathrm{2D}(\mathbf{x})}{k_\mathrm{eff}^\mathrm{3D}(\mathbf{x})}
\label{eq:lax-def}
\end{equation}

\noindent where the symbols have the following meaning.

\begin{itemize}
  \item $L_\mathrm{ax}$ is the axial leakage factor, dimensionless.

  \item $k_\mathrm{eff}^\mathrm{2D}$ is the eigenvalue of the
        two-dimensional core with the campaign radial reflector,
        dimensionless.

  \item $k_\mathrm{eff}^\mathrm{3D}$ is the eigenvalue of the same core
        at finite height with the axial reflectors, dimensionless.
\end{itemize}

Table~\ref{tab:lax-study} reports the six Campaign 6 designs studied.
They span an unrodded eigenvalue from 1.071 to 1.279 and enrichments
from 3 to \SI{14}{wt\%}, and their two-dimensional regulating-bank
worths, \num{14839} to \SI{17815}{pcm}, lie within one standard
deviation of the Campaign 7 archive median.

\begin{table}[htbp]
  \centering
  \caption{Axial leakage study on six Campaign 6 designs. Rodded
  columns have all sixteen regulating banks inserted. Subcriticality
  is $(1 - 1/k)$ in pcm, negative when supercritical.}
  \label{tab:lax-study}
  \begin{tabular}{ccccccc}
    \toprule
    Design & $k^\mathrm{2D}$ & $k^\mathrm{3D}$ & $L_\mathrm{ax}$ & $k^\mathrm{3D}_\mathrm{ARI}$ & 2D subcrit. [pcm] & 3D subcrit. [pcm] \\
    \midrule
     1 & 1.11990 & 1.08888 & 1.0285 & 0.93673 &  3821 &  6755 \\
    22 & 1.09768 & 1.06638 & 1.0293 & 0.89620 &  8534 & 11583 \\
    31 & 1.10226 & 1.07207 & 1.0282 & 0.91742 &  6095 &  9001 \\
    59 & 1.07088 & 1.04018 & 1.0295 & 0.89711 &  8402 & 11469 \\
    69 & 1.12079 & 1.08839 & 1.0298 & 0.92475 &  5095 &  8137 \\
    86 & 1.27869 & 1.24367 & 1.0282 & 1.07050 & $-9136$ & $-6585$ \\
    \bottomrule
  \end{tabular}
\end{table}

Three results carry into the method.

\begin{enumerate}
  \item The factor is effectively a constant. It ranges from
        \num{1.0282} to \num{1.0298} over a factor of two in excess
        reactivity, a spread of \SI{160}{pcm}. One value, the mean
        \num{1.0289}, therefore serves the whole design space with a
        stated systematic of $\pm\SI{80}{pcm}$, about $\pm 13$ EFPD at
        the end-of-cycle slopes of Section~\ref{sec:res-c7-uncertainty}.
        This is the same structural property that justifies Route B
        itself, namely that a leakage factor is governed by the
        geometry, and the geometry is fixed.

  \item The rodded core is \num{2900} to \SI{3100}{pcm} more
        subcritical in three dimensions than in two. The
        two-dimensional control screen of Section~\ref{sec:meth-ctrl}
        is therefore conservative by that amount, and the direction of
        the conservatism is measured rather than argued.

  \item Design 86, the high-enrichment Campaign 6 champion, is
        supercritical with all regulating banks inserted in both
        models. The long-cycle end of the unconstrained design space
        cannot be held down by the regulating system, which is the
        physical reason for the enrichment cap and the control screen
        of the final campaign.
\end{enumerate}

From Campaign 8 the axial factor multiplies the target, so that the
end of cycle is defined on a finite-height core from the start.

\begin{equation}
k_\mathrm{target}^\mathrm{3D}(t_\mathrm{refl}) = k_\mathrm{target}^\mathrm{2D}(t_\mathrm{refl}) \times L_\mathrm{ax}
\label{eq:ktarget-3d}
\end{equation}

\noindent where the symbols have the following meaning.

\begin{itemize}
  \item $k_\mathrm{target}^\mathrm{3D}$ is the target that the depleting
        assembly must retain at end of cycle for the finite core to be
        critical, dimensionless.

  \item $k_\mathrm{target}^\mathrm{2D}$ is the Route B target of
        Equation~\eqref{eq:ktarget} at the fixed pitch, dimensionless.

  \item $L_\mathrm{ax}$ is the axial factor, \num{1.0289}.

  \item $t_\mathrm{refl}$ is the reflector thickness, in cm.
\end{itemize}

With the pitch fixed at \SI{1.26}{cm} the calibration table of
Section~\ref{sec:meth-routeB} collapses to one dimension. It was
recomputed at seven reflector nodes from 2.0 to \SI{5.66}{cm}, one
assembly solve and seven core solves. The per-node statistical error
from the core solve alone is 43 to \SI{60}{pcm}, and the single
assembly solve shifts every node together by \SI{90}{pcm} without
affecting the shape. A weighted straight line describes the seven
nodes with a chi-square per degree of freedom of 0.38 and a worst
residual of 1.16 standard deviations, and the fitted values replace the
raw ones to remove a one-sigma fluctuation at \SI{3.8}{cm}.

\begin{equation}
k_\mathrm{target}^\mathrm{2D}(t_\mathrm{refl}) = 1.054528 - 0.000516\, t_\mathrm{refl}
\label{eq:ktarget-c8-fit}
\end{equation}

\noindent where the symbols have the following meaning.

\begin{itemize}
  \item $k_\mathrm{target}^\mathrm{2D}$ is the two-dimensional target at
        pitch \SI{1.26}{cm}, dimensionless.

  \item $t_\mathrm{refl}$ is the reflector thickness, in cm, over 2.0
        to 5.66.

  \item The slope is \SI{51.6}{pcm} per centimetre of reflector.
\end{itemize}

The three-dimensional target runs from 1.0839 at \SI{2.0}{cm} to
1.0820 at \SI{5.66}{cm}. Across the entire Campaign 8 design space the
target therefore moves by about \SI{200}{pcm}, roughly 0.3 MWd/kgHM or
34 EFPD, so the reflector thickness has ceased to be a reactivity
variable in this campaign and acts on the radial power shape alone.

The correction raises the target by about \SI{2900}{pcm} relative to
the two-dimensional campaigns. At the median end-of-cycle slope of
\SI{607}{pcm} per MWd/kgHM that is 4.8 MWd/kgHM, roughly 450 to 500
EFPD. Campaign 8 cycle lengths are therefore not comparable with those
of Campaigns 1 to 7 without adding that amount back, and they are the
first in this work that refer to a finite-height core.
%% <<< M6

%% ===================================================================
%% BLOCK M7  (methodology.tex, insert after line 1285, the end of
%% section "Evaluator-level radial zoning and the enrichment search box")
%% ===================================================================
%% >>> M7
\section{The Campaign 8 design space}
\label{sec:meth-c8-space}

The final campaign fixes what the earlier campaigns established and
searches only what remains open. Four decisions define it.

\begin{enumerate}
  \item The pin pitch is fixed at \SI{1.26}{cm}, the Westinghouse
        $17\times17$ value that the reference fuel design of
        Section~\ref{sec:design-basis} already adopts. Two reasons
        compel it. The guide-tube outer radius of \SI{0.6121}{cm}
        exceeds the half-pitch below \SI{1.2242}{cm}, and an OpenMC
        lattice clips a universe at its cell boundary silently, so
        every earlier design below that pitch carried guide tubes with
        up to \SI{27}{\%} of their wall replaced by moderator. And the
        pitch was the strongest single lever on the regulating-bank
        worth (Section~\ref{sec:res-c7-worth}), so fixing it removes a
        source of variability that the control screen would otherwise
        have to absorb.

  \item The two assembly enrichment zones collapse to one variable.
        The core-level enrichment variation comes entirely from the
        C, M, P ring multipliers of Section~\ref{sec:meth-zoning},
        \num{0.720}, \num{0.893} and \num{1.150}, and the cap on the
        as-built peripheral ring is raised to \SI{16}{wt\%}, so the
        search box is 3.0 to \SI{13.913}{wt\%}.

  \item The vessel-fit constraint reserves a clearance for a core
        barrel and a downcomer. The barrel is the \SI{5.08}{cm}
        stainless steel wall of the NuScale-like benchmark and the
        downcomer is \SI{2.0}{cm}, chosen as a minimum coolant return
        path, with a sensitivity to \SI{3.7}{cm} run on the candidate
        designs only. The clearance enters the geometric constraint
        and not the transport model, whose outer boundary remains a
        vacuum at the reflector surface. The barrel and the downcomer
        are represented physically in the three-dimensional
        confirmation of the candidates.

  \item The reflector thickness is re-bounded to the budget that
        remains.
\end{enumerate}

\begin{equation}
t_\mathrm{refl}^{\max} = R_\mathrm{v} - c_\mathrm{v} - R_\mathrm{env}(p) - \delta
\label{eq:refl-budget-c8}
\end{equation}

\noindent where the symbols have the following meaning.

\begin{itemize}
  \item $t_\mathrm{refl}^{\max}$ is the largest reflector thickness that
        fits, \SI{5.669}{cm}, bounded in the design space at
        \SI{5.66}{cm} so that the box edge is strictly feasible.

  \item $R_\mathrm{v}$ is the vessel internal radius, \SI{90}{cm}.

  \item $c_\mathrm{v}$ is the reserved clearance, \SI{7.08}{cm}.

  \item $R_\mathrm{env}(p)$ is the circumscribed radius of the
        32-assembly fuel envelope at pitch $p$, \SI{77.231}{cm} at
        \SI{1.26}{cm}.

  \item $\delta$ is the lattice pad, \SI{0.02}{cm}.
\end{itemize}

The consequence is stated rather than discovered. Five of the six
designs of the Campaign 7 control-screen front, and 14 of its 60
designs, are unbuildable at the fixed pitch, because the low-peaking
end of that front lives at thick reflectors that fit only at narrow
pitch. Campaign 8 therefore trades the lowest peaking values of
Campaign 7 for a geometry that is consistent, and the peaking floor is
expected to rise to about 1.60.

Two further settings close the specification. The acquisition runs
with the minimum separation of 0.14 and the feasibility margin of 1.5
standard deviations established in Campaign 6, stated explicitly on
the command line after the Campaign 7 regression. And the design of
experiments has 36 points for four variables, nine per variable,
followed by four infill blocks of six, 60 evaluations in all.
%% <<< M7

%% ===================================================================
%% BLOCK M8  (methodology.tex, tab:campaigns, insert after the C6 row,
%% line 1315)
%% ===================================================================
%% >>> M8
    C7 & 60 & Measured controllability screen $g_\mathrm{ctrl}$
       & Screens nested, derived ceiling discards controllable designs, two acquisition collapses \\
    C8 & 60 & Fixed pitch, single enrichment, axial-corrected target, two-level screen
       & Candidate front for the LABGENE-class core \\
%% <<< M8

%% ===================================================================
%% BLOCK M9  (methodology.tex, insert after line 1539, the end of
%% section "Diagnostics, environments, reproducibility")
%% ===================================================================
%% >>> M9
\subsection{Cost accounting}
\label{sec:meth-cost}

The per-evaluation cost stored in the archive is the sum of the
assembly solve, the core solve and the depletion timers. Through
Campaign 7 the control solve ran after those timers closed and was
stored separately, so the archived cost understated the true cost by
the control share, \SI{11}{\%} over Campaign 7. The block-level wall
time of the run log, which is measured around the whole evaluation
call, was unaffected and is the figure reported in the cost tables.
From Campaign 8 every rodded solve is folded into the archived cost
and the per-case log line reports the full figure.

The per-case parser of the run log was also found never to have
matched a case line of any campaign, because its regular expression
did not anticipate the censoring marker inserted after the cycle
length. The failure was silent. No quantity in this thesis depends on
its output, since every per-design value is read from the archive, and
the parser was corrected before Campaign 8.
%% <<< M9
